\documentclass[11pt,a4paper]{article}

%====================================================
% PACKAGES
%====================================================
\usepackage[utf8]{inputenc}
\usepackage[T1]{fontenc}
\usepackage[french]{babel}
\usepackage[a4paper,margin=1.8cm]{geometry}
\usepackage{amsmath,amssymb,amsthm,mathtools}
\usepackage{lmodern}
\usepackage{xcolor}
\usepackage{fancyhdr}
\usepackage{lastpage}
\usepackage{enumitem}
\usepackage{array,tabularx,multicol}
\usepackage{tikz}
\usepackage{pgfplots}
\usepackage[most]{tcolorbox}
\usepackage{hyperref}
\pgfplotsset{compat=1.18}

%====================================================
% COULEURS
%====================================================
\definecolor{bleu}{RGB}{30,80,160}
\definecolor{vert}{RGB}{20,120,80}
\definecolor{rouge}{RGB}{180,50,50}
\definecolor{orange}{RGB}{210,120,20}
\definecolor{gris}{RGB}{245,245,245}

%====================================================
% MISE EN PAGE
%====================================================
\setlength{\headheight}{14pt}
\pagestyle{fancy}
\fancyhf{}
\fancyhead[L]{Géométrie collège}
\fancyhead[C]{Cours complet}
\fancyhead[R]{Page \thepage/\pageref{LastPage}}
\fancyfoot[C]{\small Support professeur -- De la 6e à la 3e}

\hypersetup{
    colorlinks=true,
    linkcolor=bleu,
    urlcolor=bleu
}

\setlist[itemize]{label=--,leftmargin=1.2cm}
\setlist[enumerate]{leftmargin=1.2cm}
\renewcommand{\arraystretch}{1.25}

%====================================================
% ENVIRONNEMENTS
%====================================================
\newtcolorbox{definitionbox}{
  colback=blue!4!white,
  colframe=bleu,
  title=Définition,
  fonttitle=\bfseries,
  breakable
}

\newtcolorbox{proprietebox}{
  colback=green!4!white,
  colframe=vert,
  title=Propriété / Théorème,
  fonttitle=\bfseries,
  breakable
}

\newtcolorbox{methodebox}{
  colback=orange!6!white,
  colframe=orange,
  title=Méthode,
  fonttitle=\bfseries,
  breakable
}

\newtcolorbox{exemplebox}{
  colback=gray!6!white,
  colframe=black!50,
  title=Exemple corrigé,
  fonttitle=\bfseries,
  breakable
}

\newtcolorbox{preuvebox}{
  colback=white,
  colframe=black,
  title=Démonstration,
  fonttitle=\bfseries,
  breakable
}

\newtcolorbox{attentionbox}{
  colback=red!5!white,
  colframe=rouge,
  title=Erreur classique,
  fonttitle=\bfseries,
  breakable
}

\newtcolorbox{retenirbox}{
  colback=yellow!12!white,
  colframe=orange!80!black,
  title=À retenir,
  fonttitle=\bfseries,
  breakable
}

\newcounter{exercice}
\newenvironment{exercice}[1][]
{\refstepcounter{exercice}\par\medskip\noindent\textbf{Exercice \theexercice. #1}\par\smallskip}
{\medskip}

\newenvironment{corrige}[1][]
{\par\medskip\noindent\textbf{Corrigé de l'exercice #1.}\par\smallskip}
{\medskip}

\newcommand{\chapitre}[1]{
\clearpage
\section{#1}
}

\newcommand{\objectif}[1]{\subsection*{A. Objectifs du chapitre}#1}
\newcommand{\defs}{\subsection*{B. Définitions essentielles}}
\newcommand{\props}{\subsection*{C. Propriétés et théorèmes}}
\newcommand{\demos}{\subsection*{D. Démonstrations rédigées}}
\newcommand{\methodes}{\subsection*{E. Méthodes types}}
\newcommand{\exemples}{\subsection*{F. Exemples corrigés pas à pas}}
\newcommand{\erreurs}{\subsection*{G. Erreurs fréquentes}}
\newcommand{\exos}{\subsection*{H. Exercices d'application progressifs}}
\newcommand{\exosdemo}{\subsection*{I. Exercices de démonstration}}
\newcommand{\corriges}{\subsection*{J. Corrigés détaillés}}

%====================================================
% DOCUMENT
%====================================================
\begin{document}

\begin{titlepage}
\centering
\vspace*{2cm}
{\Huge\bfseries Cours complet de géométrie au collège\par}
\vspace{0.4cm}
{\Large De la 6e à la 3e\par}
\vspace{1cm}
{\large Définitions, propriétés, démonstrations, méthodes, exemples corrigés et exercices\par}
\vfill

\vfill
{\large Année scolaire \the\year}
\end{titlepage}

\tableofcontents
\clearpage

%====================================================
\section*{Préambule : faire de la géométrie}
\addcontentsline{toc}{section}{Préambule : faire de la géométrie}

La géométrie n'est pas seulement l'art de tracer de belles figures. Une figure permet d'observer, de conjecturer et de chercher une idée. Mais une figure ne prouve pas. Démontrer, c'est expliquer pourquoi une conclusion est vraie en utilisant des définitions, des propriétés ou des théorèmes connus.

\begin{retenirbox}
En géométrie, une rédaction solide suit souvent la forme :
\[
\text{Je sais que} \quad \longrightarrow \quad \text{J'utilise la propriété} \quad \longrightarrow \quad \text{Donc}
\]
\end{retenirbox}

\begin{exemplebox}
On observe sur une figure que deux segments semblent avoir la même longueur. Ce n'est pas une preuve. Pour prouver que $AB=AC$, on peut par exemple montrer que $A$ appartient à la médiatrice de $[BC]$, ou que le triangle $ABC$ est isocèle en $A$, ou utiliser une symétrie qui transforme $B$ en $C$.
\end{exemplebox}

%====================================================
\chapitre{Les bases de la géométrie}

\objectif{
\begin{itemize}
\item Savoir nommer les objets géométriques fondamentaux.
\item Comprendre la différence entre point, droite, demi-droite et segment.
\item Savoir coder une figure et exploiter les codages.
\item Savoir utiliser la notion de milieu.
\end{itemize}}

\defs
\begin{definitionbox}
Un \textbf{point} est un emplacement du plan. On le note avec une lettre majuscule : $A$, $B$, $C$.

Une \textbf{droite} est illimitée dans les deux sens. La droite passant par $A$ et $B$ se note $(AB)$.

Le \textbf{segment} d'extrémités $A$ et $B$ est la partie de la droite $(AB)$ comprise entre $A$ et $B$. Il se note $[AB]$.

La \textbf{demi-droite} d'origine $A$ passant par $B$ se note $[AB)$.
\end{definitionbox}

\begin{definitionbox}
Trois points $A$, $B$ et $C$ sont \textbf{alignés} lorsqu'ils appartiennent à une même droite.

Un point $M$ est le \textbf{milieu} du segment $[AB]$ lorsque :
\[
M \in [AB] \quad \text{et} \quad AM=MB.
\]
\end{definitionbox}

\props
\begin{proprietebox}
Par deux points distincts, il passe une unique droite.
\end{proprietebox}

\begin{proprietebox}
Si $M$ est le milieu de $[AB]$, alors :
\[
AM=MB=\frac{AB}{2}.
\]
Réciproquement, si $M\in[AB]$ et $AM=MB$, alors $M$ est le milieu de $[AB]$.
\end{proprietebox}

\demos
\begin{preuvebox}
La propriété du milieu vient directement de la définition. Dire que $M$ est le milieu de $[AB]$ signifie deux choses : d'abord $M$ est sur le segment $[AB]$, ensuite il partage ce segment en deux segments de même longueur. Donc $AM=MB$. Comme $AB=AM+MB$, on obtient $AB=2AM$, donc $AM=\frac{AB}{2}$.
\end{preuvebox}

\methodes
\begin{methodebox}
Pour prouver qu'un point $M$ est le milieu de $[AB]$, il faut prouver deux informations :
\begin{enumerate}
\item $M$ appartient au segment $[AB]$ ;
\item $AM=MB$.
\end{enumerate}
Une seule de ces deux informations ne suffit pas.
\end{methodebox}

\exemples
\begin{exemplebox}
On sait que $AB=8$ cm et que $M$ est le milieu de $[AB]$. Alors :
\[
AM=MB=\frac{AB}{2}=\frac{8}{2}=4 \text{ cm}.
\]
\end{exemplebox}

\erreurs
\begin{attentionbox}
Dire que $AM=MB$ ne suffit pas toujours à prouver que $M$ est le milieu de $[AB]$. Il faut aussi vérifier que $M$ appartient au segment $[AB]$.
\end{attentionbox}

\exos
\begin{exercice}[Milieu]
On donne un segment $[AB]$ de longueur $12$ cm. Le point $M$ est le milieu de $[AB]$.
\begin{enumerate}
\item Calculer $AM$ et $MB$.
\item Rédiger une phrase justifiant le calcul.
\end{enumerate}
\end{exercice}

\begin{exercice}[Alignement]
Les points $A$, $B$ et $C$ sont alignés dans cet ordre. On sait que $AB=5$ cm et $BC=7$ cm. Calculer $AC$.
\end{exercice}

\exosdemo
\begin{exercice}[Démontrer un milieu]
Les points $A$, $M$ et $B$ sont alignés dans cet ordre. On sait que $AM=6$ cm et $AB=12$ cm. Démontrer que $M$ est le milieu de $[AB]$.
\end{exercice}

\corriges
\begin{corrige}[1]
Puisque $M$ est le milieu de $[AB]$, on a $AM=MB=\dfrac{AB}{2}$. Or $AB=12$ cm, donc $AM=MB=6$ cm.
\end{corrige}

\begin{corrige}[2]
Les points $A$, $B$ et $C$ sont alignés dans cet ordre, donc $AC=AB+BC$. Ainsi $AC=5+7=12$ cm.
\end{corrige}

\begin{corrige}[3]
On sait que $A$, $M$ et $B$ sont alignés dans cet ordre, donc $M\in[AB]$. De plus $AB=12$ cm et $AM=6$ cm, donc $MB=AB-AM=12-6=6$ cm. Ainsi $AM=MB$ et $M\in[AB]$. Donc $M$ est le milieu de $[AB]$.
\end{corrige}

%====================================================
\chapitre{Les angles}

\objectif{
\begin{itemize}
\item Savoir mesurer et nommer un angle.
\item Reconnaître les angles particuliers.
\item Utiliser les propriétés liées aux droites parallèles.
\item Rédiger des démonstrations avec les angles.
\end{itemize}}

\defs
\begin{definitionbox}
Un \textbf{angle} est une portion de plan délimitée par deux demi-droites de même origine. L'origine commune est le \textbf{sommet} de l'angle.

L'angle de sommet $B$ formé par les demi-droites $[BA)$ et $[BC)$ se note $\widehat{ABC}$.
\end{definitionbox}

\begin{definitionbox}
Un angle est :
\begin{itemize}
\item \textbf{aigu} si sa mesure est comprise entre $0^\circ$ et $90^\circ$ ;
\item \textbf{droit} si sa mesure vaut $90^\circ$ ;
\item \textbf{obtus} si sa mesure est comprise entre $90^\circ$ et $180^\circ$ ;
\item \textbf{plat} si sa mesure vaut $180^\circ$.
\end{itemize}
\end{definitionbox}

\props
\begin{proprietebox}
Deux angles opposés par le sommet ont la même mesure.
\end{proprietebox}

\begin{proprietebox}
Si deux droites parallèles sont coupées par une sécante, alors les angles alternes-internes sont égaux, et les angles correspondants sont égaux.
\end{proprietebox}

\begin{proprietebox}
Réciproquement, si deux angles alternes-internes formés par deux droites et une sécante sont égaux, alors les deux droites sont parallèles.
\end{proprietebox}

\demos
\begin{preuvebox}
Considérons deux angles opposés par le sommet. Chacun de ces angles complète le même angle adjacent pour former un angle plat de $180^\circ$. Si on note les deux angles opposés $x$ et $y$, et l'angle adjacent $z$, on a :
\[
x+z=180^\circ \quad \text{et} \quad y+z=180^\circ.
\]
Donc $x+z=y+z$, d'où $x=y$.
\end{preuvebox}

\methodes
\begin{methodebox}
Pour démontrer que deux droites sont parallèles avec les angles :
\begin{enumerate}
\item repérer une sécante commune ;
\item identifier deux angles alternes-internes ou correspondants ;
\item montrer que ces angles ont la même mesure ;
\item conclure que les droites sont parallèles.
\end{enumerate}
\end{methodebox}

\exemples
\begin{exemplebox}
Les droites $(d)$ et $(d')$ sont coupées par une sécante. Deux angles alternes-internes mesurent chacun $62^\circ$. Comme deux angles alternes-internes sont égaux, les droites $(d)$ et $(d')$ sont parallèles.
\end{exemplebox}

\erreurs
\begin{attentionbox}
Deux angles qui se ressemblent sur une figure ne sont pas forcément égaux. Il faut une propriété : angles opposés par le sommet, angles correspondants, angles alternes-internes avec parallèles, etc.
\end{attentionbox}

\exos
\begin{exercice}[Angles supplémentaires]
Deux angles sont supplémentaires. Le premier mesure $118^\circ$. Calculer la mesure du second.
\end{exercice}

\begin{exercice}[Angles opposés]
Deux droites se coupent en $O$. Un angle mesure $47^\circ$. Quelle est la mesure de l'angle opposé par le sommet ? Justifier.
\end{exercice}

\exosdemo
\begin{exercice}[Parallélisme]
Deux droites $(d)$ et $(d')$ sont coupées par une sécante. Deux angles correspondants mesurent $74^\circ$. Démontrer que $(d)$ et $(d')$ sont parallèles.
\end{exercice}

\corriges
\begin{corrige}[4]
Deux angles supplémentaires ont une somme égale à $180^\circ$. Le second angle mesure donc $180^\circ-118^\circ=62^\circ$.
\end{corrige}

\begin{corrige}[5]
Deux angles opposés par le sommet ont la même mesure. L'angle opposé mesure donc $47^\circ$.
\end{corrige}

\begin{corrige}[6]
Les deux angles correspondants formés par les deux droites et la sécante sont égaux. Or si deux angles correspondants sont égaux, alors les deux droites sont parallèles. Donc $(d)\parallel(d')$.
\end{corrige}

%====================================================
\chapitre{Droites perpendiculaires et parallèles}

\objectif{
\begin{itemize}
\item Reconnaître et construire des droites parallèles ou perpendiculaires.
\item Utiliser les propriétés fondamentales de perpendicularité et de parallélisme.
\item Rédiger une preuve simple avec ces propriétés.
\end{itemize}}

\defs
\begin{definitionbox}
Deux droites sont \textbf{perpendiculaires} lorsqu'elles se coupent en formant un angle droit.

Deux droites sont \textbf{parallèles} lorsqu'elles ne se coupent jamais, même si on les prolonge indéfiniment.
\end{definitionbox}

\props
\begin{proprietebox}
Si deux droites sont perpendiculaires à une même troisième droite, alors elles sont parallèles.
\end{proprietebox}

\begin{proprietebox}
Si deux droites sont parallèles, alors toute droite perpendiculaire à l'une est perpendiculaire à l'autre.
\end{proprietebox}

\begin{proprietebox}
Si deux droites sont parallèles à une même troisième droite, alors elles sont parallèles entre elles.
\end{proprietebox}

\demos
\begin{preuvebox}
La première propriété repose sur l'idée suivante : deux droites perpendiculaires à une même droite forment chacune un angle droit avec cette droite. Elles ont donc la même direction. Deux droites distinctes ayant la même direction sont parallèles.
\end{preuvebox}

\methodes
\begin{methodebox}
Pour démontrer que deux droites sont parallèles, on peut chercher :
\begin{itemize}
\item deux droites perpendiculaires à une même droite ;
\item deux droites parallèles à une même droite ;
\item des angles alternes-internes ou correspondants égaux ;
\item une configuration de Thalès réciproque.
\end{itemize}
\end{methodebox}

\exemples
\begin{exemplebox}
On sait que $(d_1)\perp(\Delta)$ et $(d_2)\perp(\Delta)$. Alors $(d_1)$ et $(d_2)$ sont toutes deux perpendiculaires à la même droite $(\Delta)$. Donc $(d_1)\parallel(d_2)$.
\end{exemplebox}

\erreurs
\begin{attentionbox}
Ne pas confondre : deux droites perpendiculaires à une même droite sont parallèles, mais deux droites perpendiculaires entre elles ne sont pas parallèles.
\end{attentionbox}

\exos
\begin{exercice}[Droites parallèles]
On sait que $(d_1)\perp(\Delta)$ et $(d_2)\perp(\Delta)$. Que peut-on dire de $(d_1)$ et $(d_2)$ ? Justifier.
\end{exercice}

\exosdemo
\begin{exercice}[Preuve de perpendicularité]
On sait que $(d)\parallel(d')$ et que $(\Delta)\perp(d)$. Démontrer que $(\Delta)\perp(d')$.
\end{exercice}

\corriges
\begin{corrige}[7]
Si deux droites sont perpendiculaires à une même troisième droite, alors elles sont parallèles. Donc $(d_1)\parallel(d_2)$.
\end{corrige}

\begin{corrige}[8]
On sait que $(d)\parallel(d')$ et que $(\Delta)\perp(d)$. Or si deux droites sont parallèles, alors toute droite perpendiculaire à l'une est perpendiculaire à l'autre. Donc $(\Delta)\perp(d')$.
\end{corrige}

%====================================================
\chapitre{Triangles}

\objectif{
\begin{itemize}
\item Connaître les triangles particuliers et leurs propriétés.
\item Utiliser la somme des angles d'un triangle.
\item Construire et raisonner avec les droites remarquables.
\item Démontrer des propriétés dans les triangles.
\end{itemize}}

\defs
\begin{definitionbox}
Un \textbf{triangle} est un polygone à trois côtés. Il possède trois sommets, trois côtés et trois angles.

Un triangle est :
\begin{itemize}
\item \textbf{isocèle} s'il possède deux côtés de même longueur ;
\item \textbf{équilatéral} s'il possède trois côtés de même longueur ;
\item \textbf{rectangle} s'il possède un angle droit.
\end{itemize}
\end{definitionbox}

\begin{definitionbox}
Dans un triangle :
\begin{itemize}
\item une \textbf{hauteur} est une droite passant par un sommet et perpendiculaire au côté opposé ;
\item une \textbf{médiane} est une droite passant par un sommet et par le milieu du côté opposé ;
\item une \textbf{médiatrice} est une droite perpendiculaire à un côté en son milieu ;
\item une \textbf{bissectrice} partage un angle en deux angles de même mesure.
\end{itemize}
\end{definitionbox}

\props
\begin{proprietebox}
Dans tout triangle, la somme des mesures des trois angles vaut $180^\circ$.
\end{proprietebox}

\begin{proprietebox}
Dans un triangle isocèle, les angles à la base sont égaux.

Réciproquement, si un triangle possède deux angles égaux, alors il est isocèle.
\end{proprietebox}

\begin{proprietebox}
Dans un triangle équilatéral, les trois angles mesurent $60^\circ$.
\end{proprietebox}

\begin{proprietebox}
Dans un triangle, la longueur d'un côté est toujours strictement inférieure à la somme des longueurs des deux autres côtés. C'est l'inégalité triangulaire.
\end{proprietebox}

\demos
\begin{preuvebox}[title=Démonstration de la somme des angles]
Soit $ABC$ un triangle. On trace par $A$ la droite parallèle à $(BC)$. Les angles formés avec les côtés $[AB]$ et $[AC]$ sont égaux aux angles en $B$ et en $C$ du triangle, car ce sont des angles alternes-internes formés par des droites parallèles. Ces trois angles, placés autour de $A$ sur une même droite, forment un angle plat. Leur somme vaut donc $180^\circ$.
\end{preuvebox}

\methodes
\begin{methodebox}
Pour calculer un angle dans un triangle :
\begin{enumerate}
\item identifier les angles connus ;
\item utiliser éventuellement les propriétés des triangles particuliers ;
\item appliquer la somme des angles :
\[
\widehat{A}+\widehat{B}+\widehat{C}=180^\circ.
\]
\end{enumerate}
\end{methodebox}

\exemples
\begin{exemplebox}
Dans un triangle $ABC$, on sait que $\widehat{A}=43^\circ$ et $\widehat{B}=72^\circ$. Alors :
\[
\widehat{C}=180^\circ-43^\circ-72^\circ=65^\circ.
\]
\end{exemplebox}

\erreurs
\begin{attentionbox}
Un triangle isocèle possède deux côtés égaux et deux angles égaux, mais il faut faire attention à la base. Si $AB=AC$, alors les angles égaux sont les angles en $B$ et en $C$.
\end{attentionbox}

\exos
\begin{exercice}[Somme des angles]
Dans un triangle $ABC$, $\widehat{A}=38^\circ$ et $\widehat{B}=81^\circ$. Calculer $\widehat{C}$.
\end{exercice}

\begin{exercice}[Triangle isocèle]
Le triangle $ABC$ est isocèle en $A$. On sait que $\widehat{B}=54^\circ$. Calculer les deux autres angles.
\end{exercice}

\exosdemo
\begin{exercice}[Reconnaissance d'un triangle isocèle]
Dans un triangle $ABC$, on sait que $\widehat{B}=\widehat{C}$. Démontrer que $ABC$ est isocèle.
\end{exercice}

\corriges
\begin{corrige}[9]
Dans un triangle, la somme des angles vaut $180^\circ$. Donc $\widehat{C}=180^\circ-38^\circ-81^\circ=61^\circ$.
\end{corrige}

\begin{corrige}[10]
Le triangle est isocèle en $A$, donc $AB=AC$ et les angles à la base sont égaux : $\widehat{B}=\widehat{C}$. Ainsi $\widehat{C}=54^\circ$. Puis $\widehat{A}=180^\circ-54^\circ-54^\circ=72^\circ$.
\end{corrige}

\begin{corrige}[11]
Si un triangle possède deux angles égaux, alors il est isocèle. Comme $\widehat{B}=\widehat{C}$, le triangle $ABC$ est isocèle en $A$.
\end{corrige}

%====================================================
\chapitre{Quadrilatères}

\objectif{
\begin{itemize}
\item Identifier les quadrilatères usuels.
\item Connaître les propriétés des côtés, angles et diagonales.
\item Utiliser les conditions suffisantes de reconnaissance.
\end{itemize}}

\defs
\begin{definitionbox}
Un \textbf{quadrilatère} est un polygone à quatre côtés.

Un \textbf{parallélogramme} est un quadrilatère dont les côtés opposés sont parallèles deux à deux.

Un \textbf{rectangle} est un quadrilatère ayant quatre angles droits.

Un \textbf{losange} est un quadrilatère ayant quatre côtés de même longueur.

Un \textbf{carré} est à la fois un rectangle et un losange.
\end{definitionbox}

\props
\begin{proprietebox}
Dans un parallélogramme :
\begin{itemize}
\item les côtés opposés sont parallèles ;
\item les côtés opposés ont la même longueur ;
\item les diagonales se coupent en leur milieu ;
\item les angles opposés sont égaux.
\end{itemize}
\end{proprietebox}

\begin{proprietebox}
Si un quadrilatère a ses diagonales qui se coupent en leur milieu, alors c'est un parallélogramme.
\end{proprietebox}

\begin{proprietebox}
Un parallélogramme ayant un angle droit est un rectangle.

Un parallélogramme ayant deux côtés consécutifs de même longueur est un losange.

Un quadrilatère qui est à la fois rectangle et losange est un carré.
\end{proprietebox}

\demos
\begin{preuvebox}
Dans un parallélogramme, les côtés opposés sont parallèles par définition. Les propriétés des diagonales et des longueurs découlent du fait que la translation qui envoie un sommet sur le sommet voisin conserve les longueurs, le parallélisme et les milieux. Au collège, on les utilise comme propriétés admises ou vérifiées par construction.
\end{preuvebox}

\methodes
\begin{methodebox}
Pour prouver qu'un quadrilatère est un parallélogramme, on peut montrer :
\begin{itemize}
\item que ses côtés opposés sont parallèles deux à deux ;
\item que ses diagonales se coupent en leur milieu ;
\item que ses côtés opposés ont la même longueur deux à deux.
\end{itemize}
\end{methodebox}

\exemples
\begin{exemplebox}
Dans un quadrilatère $ABCD$, les diagonales $[AC]$ et $[BD]$ se coupent en $O$. On sait que $OA=OC$ et $OB=OD$. Donc les diagonales se coupent en leur milieu. Par conséquent, $ABCD$ est un parallélogramme.
\end{exemplebox}

\erreurs
\begin{attentionbox}
Un carré est toujours un rectangle, mais un rectangle n'est pas toujours un carré. Un carré est toujours un losange, mais un losange n'est pas toujours un carré.
\end{attentionbox}

\exos
\begin{exercice}[Parallélogramme]
Dans un quadrilatère $ABCD$, les diagonales se coupent en $O$. On sait que $OA=OC$ et $OB=OD$. Que peut-on conclure ?
\end{exercice}

\exosdemo
\begin{exercice}[Reconnaissance d'un rectangle]
$ABCD$ est un parallélogramme et $\widehat{ABC}=90^\circ$. Démontrer que $ABCD$ est un rectangle.
\end{exercice}

\corriges
\begin{corrige}[12]
Puisque $OA=OC$, $O$ est le milieu de $[AC]$. Puisque $OB=OD$, $O$ est le milieu de $[BD]$. Les diagonales se coupent donc en leur milieu. Ainsi $ABCD$ est un parallélogramme.
\end{corrige}

\begin{corrige}[13]
On sait que $ABCD$ est un parallélogramme et qu'il possède un angle droit. Or un parallélogramme qui possède un angle droit est un rectangle. Donc $ABCD$ est un rectangle.
\end{corrige}

%====================================================
\chapitre{Cercles et médiatrices}

\objectif{
\begin{itemize}
\item Utiliser le vocabulaire du cercle.
\item Comprendre la propriété caractéristique de la médiatrice.
\item Construire un cercle circonscrit.
\item Raisonner avec des distances égales.
\end{itemize}}

\defs
\begin{definitionbox}
Un \textbf{cercle} de centre $O$ et de rayon $r$ est l'ensemble des points situés à la distance $r$ de $O$.

Un \textbf{rayon} est un segment reliant le centre à un point du cercle.

Un \textbf{diamètre} est une corde passant par le centre. Sa longueur vaut deux fois le rayon.
\end{definitionbox}

\begin{definitionbox}
La \textbf{médiatrice} d'un segment est la droite perpendiculaire à ce segment en son milieu.
\end{definitionbox}

\props
\begin{proprietebox}
Si un point appartient à la médiatrice de $[AB]$, alors il est équidistant de $A$ et de $B$.
\[
M\in \text{médiatrice de }[AB] \Longrightarrow MA=MB.
\]
Réciproquement, si $MA=MB$, alors $M$ appartient à la médiatrice de $[AB]$.
\end{proprietebox}

\begin{proprietebox}
Les médiatrices des trois côtés d'un triangle sont concourantes. Leur point d'intersection est le centre du cercle circonscrit au triangle.
\end{proprietebox}

\demos
\begin{preuvebox}
Soit $d$ la médiatrice de $[AB]$ et $M$ un point de $d$. La droite $d$ est perpendiculaire à $[AB]$ en son milieu $I$. Les triangles $AIM$ et $BIM$ sont rectangles en $I$, ont le côté $IM$ commun et vérifient $IA=IB$. Ils sont donc superposables. Par conséquent $MA=MB$.
\end{preuvebox}

\methodes
\begin{methodebox}
Pour prouver qu'un point est sur la médiatrice de $[AB]$, on peut montrer qu'il est à égale distance de $A$ et de $B$.
\end{methodebox}

\exemples
\begin{exemplebox}
Si $PA=PB$, alors $P$ est équidistant de $A$ et de $B$. Donc $P$ appartient à la médiatrice de $[AB]$.
\end{exemplebox}

\erreurs
\begin{attentionbox}
La médiatrice d'un segment n'est pas une médiane de triangle. Une médiatrice est toujours perpendiculaire au segment et passe par son milieu.
\end{attentionbox}

\exos
\begin{exercice}[Cercle]
Un cercle a pour rayon $4$ cm. Quelle est la longueur de son diamètre ?
\end{exercice}

\exosdemo
\begin{exercice}[Médiatrice]
On sait que $PA=PB$. Démontrer que $P$ appartient à la médiatrice de $[AB]$.
\end{exercice}

\corriges
\begin{corrige}[14]
Le diamètre vaut deux fois le rayon. Donc $d=2\times 4=8$ cm.
\end{corrige}

\begin{corrige}[15]
On sait que $PA=PB$. Donc $P$ est équidistant de $A$ et de $B$. Or l'ensemble des points équidistants de $A$ et de $B$ est la médiatrice de $[AB]$. Donc $P$ appartient à la médiatrice de $[AB]$.
\end{corrige}

%====================================================
\chapitre{Symétries et transformations}

\objectif{
\begin{itemize}
\item Reconnaître les principales transformations du collège.
\item Comprendre ce qui est conservé par une transformation.
\item Utiliser les transformations pour démontrer.
\end{itemize}}

\defs
\begin{definitionbox}
La \textbf{symétrie axiale} d'axe $(d)$ transforme un point $A$ en un point $A'$ tel que $(d)$ soit la médiatrice de $[AA']$.

La \textbf{symétrie centrale} de centre $O$ transforme un point $A$ en un point $A'$ tel que $O$ soit le milieu de $[AA']$.
\end{definitionbox}

\begin{definitionbox}
Une \textbf{translation} fait glisser une figure selon une direction, un sens et une longueur.

Une \textbf{rotation} fait tourner une figure autour d'un centre selon un angle donné.

Une \textbf{homothétie} de centre $O$ et de rapport $k$ agrandit ou réduit les longueurs dans le rapport $|k|$.
\end{definitionbox}

\props
\begin{proprietebox}
Les symétries, translations et rotations conservent les longueurs, les angles, l'alignement, les milieux, le parallélisme et les aires.
\end{proprietebox}

\begin{proprietebox}
Une homothétie de rapport $k$ multiplie les longueurs par $|k|$ et les aires par $k^2$. Elle conserve les angles, l'alignement et le parallélisme.
\end{proprietebox}

\demos
\begin{preuvebox}
Dans une symétrie axiale, l'axe est la médiatrice de chaque segment reliant un point à son image. Ainsi les distances à l'axe sont conservées de part et d'autre. La figure image est superposable à la figure initiale par retournement : les longueurs et les angles sont donc conservés.
\end{preuvebox}

\methodes
\begin{methodebox}
Pour démontrer deux longueurs égales avec une symétrie, on peut montrer qu'un segment est l'image de l'autre par cette symétrie, puis utiliser la conservation des longueurs.
\end{methodebox}

\exemples
\begin{exemplebox}
Si $A'$ est l'image de $A$ et $B'$ l'image de $B$ par une symétrie axiale, alors le segment $[A'B']$ est l'image du segment $[AB]$. La symétrie conserve les longueurs, donc $A'B'=AB$.
\end{exemplebox}

\erreurs
\begin{attentionbox}
Une homothétie ne conserve pas les longueurs sauf si le rapport vaut $1$ ou $-1$. Elle multiplie les longueurs par la valeur absolue du rapport.
\end{attentionbox}

\exos
\begin{exercice}[Symétrie centrale]
$A'$ est l'image de $A$ par la symétrie de centre $O$. Que peut-on dire du point $O$ par rapport au segment $[AA']$ ?
\end{exercice}

\exosdemo
\begin{exercice}[Conservation des longueurs]
Par une symétrie axiale, $A$ a pour image $A'$ et $B$ a pour image $B'$. Démontrer que $AB=A'B'$.
\end{exercice}

\corriges
\begin{corrige}[16]
Dans une symétrie centrale de centre $O$, le centre $O$ est le milieu du segment reliant un point et son image. Donc $O$ est le milieu de $[AA']$.
\end{corrige}

\begin{corrige}[17]
La symétrie axiale conserve les longueurs. Comme $[A'B']$ est l'image de $[AB]$, on obtient $A'B'=AB$.
\end{corrige}

%====================================================
\chapitre{Aires et périmètres}

\objectif{
\begin{itemize}
\item Calculer des périmètres et des aires.
\item Choisir la bonne formule.
\item Décomposer une figure complexe.
\item Rédiger un calcul géométrique.
\end{itemize}}

\defs
\begin{definitionbox}
Le \textbf{périmètre} d'une figure est la longueur de son contour.

L'\textbf{aire} d'une figure mesure la surface qu'elle occupe.
\end{definitionbox}

\props
\begin{proprietebox}
Formules usuelles :
\[
\mathcal A_{rectangle}=L\times l,\quad \mathcal A_{triangle}=\frac{base\times hauteur}{2},
\]
\[
\mathcal A_{parallélogramme}=base\times hauteur,
\quad \mathcal A_{disque}=\pi r^2,
\quad P_{cercle}=2\pi r.
\]
\end{proprietebox}

\demos
\begin{preuvebox}
L'aire d'un triangle se comprend en le doublant. Deux triangles identiques peuvent former un parallélogramme de même base et de même hauteur. L'aire du parallélogramme vaut $base\times hauteur$, donc l'aire d'un des deux triangles vaut la moitié :
\[
\mathcal A=\frac{base\times hauteur}{2}.
\]
\end{preuvebox}

\methodes
\begin{methodebox}
Pour calculer l'aire d'une figure composée :
\begin{enumerate}
\item découper la figure en figures simples ;
\item calculer chaque aire séparément ;
\item additionner ou soustraire les aires selon la situation ;
\item écrire l'unité carrée.
\end{enumerate}
\end{methodebox}

\exemples
\begin{exemplebox}
Un triangle a une base de $10$ cm et une hauteur de $6$ cm. Son aire vaut :
\[
\mathcal A=\frac{10\times 6}{2}=30\text{ cm}^2.
\]
\end{exemplebox}

\erreurs
\begin{attentionbox}
La hauteur d'un triangle doit être perpendiculaire à la base. Un côté incliné n'est pas forcément une hauteur.
\end{attentionbox}

\exos
\begin{exercice}[Aire d'un triangle]
Calculer l'aire d'un triangle de base $8$ cm et de hauteur $5$ cm.
\end{exercice}

\exosdemo
\begin{exercice}[Formule du triangle]
Expliquer pourquoi l'aire d'un triangle est égale à la moitié de l'aire d'un parallélogramme de même base et de même hauteur.
\end{exercice}

\corriges
\begin{corrige}[18]
On utilise la formule $\mathcal A=\dfrac{base\times hauteur}{2}$. Donc $\mathcal A=\dfrac{8\times5}{2}=20\text{ cm}^2$.
\end{corrige}

\begin{corrige}[19]
Deux triangles superposables peuvent former un parallélogramme. Le parallélogramme obtenu a la même base et la même hauteur que le triangle. Son aire vaut $base\times hauteur$. Le triangle représentant la moitié de ce parallélogramme, son aire vaut $\dfrac{base\times hauteur}{2}$.
\end{corrige}

%====================================================
\chapitre{Volumes et géométrie dans l'espace}

\objectif{
\begin{itemize}
\item Identifier les solides usuels.
\item Calculer des volumes.
\item Comprendre les patrons et la perspective cavalière.
\item Raisonner dans l'espace.
\end{itemize}}

\defs
\begin{definitionbox}
Un \textbf{solide} est un objet géométrique en trois dimensions. Il peut avoir des faces, des arêtes et des sommets.

Le cube, le pavé droit, le prisme droit, le cylindre, la pyramide, le cône et la boule sont des solides usuels.
\end{definitionbox}

\props
\begin{proprietebox}
Formules usuelles de volume :
\[
V_{pavé}=L\times l\times h,
\quad V_{prisme}=\mathcal A_{base}\times h,
\quad V_{cylindre}=\pi r^2h,
\]
\[
V_{pyramide}=\frac{\mathcal A_{base}\times h}{3},
\quad V_{cône}=\frac{\pi r^2h}{3},
\quad V_{boule}=\frac{4}{3}\pi r^3.
\]
\end{proprietebox}

\methodes
\begin{methodebox}
Pour calculer le volume d'un prisme ou d'un cylindre :
\begin{enumerate}
\item calculer l'aire de la base ;
\item multiplier par la hauteur ;
\item écrire une unité cubique.
\end{enumerate}
\end{methodebox}

\exemples
\begin{exemplebox}
Un cylindre a un rayon de $3$ cm et une hauteur de $10$ cm. Son volume vaut :
\[
V=\pi r^2h=\pi\times3^2\times10=90\pi\text{ cm}^3.
\]
\end{exemplebox}

\erreurs
\begin{attentionbox}
Les unités de volume sont des unités cubiques : cm$^3$, m$^3$, mm$^3$. Une aire s'exprime en unités carrées, un volume en unités cubiques.
\end{attentionbox}

\exos
\begin{exercice}[Volume d'un pavé]
Calculer le volume d'un pavé droit de longueur $8$ cm, largeur $5$ cm et hauteur $4$ cm.
\end{exercice}

\exosdemo
\begin{exercice}[Prisme]
Expliquer pourquoi le volume d'un prisme droit est égal à l'aire de sa base multipliée par sa hauteur.
\end{exercice}

\corriges
\begin{corrige}[20]
$V=L\times l\times h=8\times5\times4=160\text{ cm}^3$.
\end{corrige}

\begin{corrige}[21]
Un prisme droit peut être vu comme un empilement de bases identiques. Si l'aire de la base est $\mathcal A_{base}$ et que la hauteur est $h$, alors le volume correspond à $h$ couches de même aire. On obtient $V=\mathcal A_{base}\times h$.
\end{corrige}

%====================================================
\chapitre{Théorème de Pythagore}

\objectif{
\begin{itemize}
\item Utiliser le théorème de Pythagore pour calculer une longueur.
\item Utiliser la réciproque pour prouver qu'un triangle est rectangle.
\item Utiliser la contraposée pour prouver qu'un triangle n'est pas rectangle.
\end{itemize}}

\defs
\begin{definitionbox}
Dans un triangle rectangle, l'\textbf{hypoténuse} est le côté opposé à l'angle droit. C'est le plus long côté du triangle.
\end{definitionbox}

\props
\begin{proprietebox}[title=Théorème de Pythagore]
Si un triangle $ABC$ est rectangle en $A$, alors :
\[
BC^2=AB^2+AC^2.
\]
\end{proprietebox}

\begin{proprietebox}[title=Réciproque]
Si, dans un triangle $ABC$, on a :
\[
BC^2=AB^2+AC^2,
\]
alors le triangle $ABC$ est rectangle en $A$.
\end{proprietebox}

\begin{proprietebox}[title=Contraposée]
Si, dans un triangle, le carré du plus long côté n'est pas égal à la somme des carrés des deux autres côtés, alors le triangle n'est pas rectangle.
\end{proprietebox}

\demos
\begin{preuvebox}
On considère un carré de côté $a+b$ découpé en quatre triangles rectangles identiques de côtés de l'angle droit $a$ et $b$, et d'hypoténuse $c$. L'aire du grand carré vaut $(a+b)^2$. Elle vaut aussi l'aire des quatre triangles plus celle du carré central de côté $c$ :
\[
(a+b)^2=4\times\frac{ab}{2}+c^2.
\]
Donc :
\[
a^2+2ab+b^2=2ab+c^2.
\]
En soustrayant $2ab$, on obtient :
\[
a^2+b^2=c^2.
\]
\end{preuvebox}

\methodes
\begin{methodebox}
Pour calculer une longueur avec Pythagore :
\begin{enumerate}
\item vérifier que le triangle est rectangle ;
\item identifier l'hypoténuse ;
\item écrire l'égalité de Pythagore ;
\item remplacer par les valeurs connues ;
\item résoudre et prendre la racine carrée si nécessaire.
\end{enumerate}
\end{methodebox}

\exemples
\begin{exemplebox}
$ABC$ est rectangle en $A$, $AB=3$ cm et $AC=4$ cm. L'hypoténuse est $BC$. Donc :
\[
BC^2=AB^2+AC^2=3^2+4^2=9+16=25.
\]
Ainsi $BC=\sqrt{25}=5$ cm.
\end{exemplebox}

\erreurs
\begin{attentionbox}
On ne peut utiliser le théorème de Pythagore que si le triangle est déjà connu comme rectangle. Pour prouver qu'il est rectangle, on utilise la réciproque.
\end{attentionbox}

\exos
\begin{exercice}[Calcul d'hypoténuse]
Un triangle $ABC$ est rectangle en $A$. On sait que $AB=6$ cm et $AC=8$ cm. Calculer $BC$.
\end{exercice}

\begin{exercice}[Réciproque]
Un triangle a pour côtés $5$ cm, $12$ cm et $13$ cm. Est-il rectangle ? Justifier.
\end{exercice}

\exosdemo
\begin{exercice}[Non-rectangle]
Un triangle a pour côtés $4$ cm, $7$ cm et $9$ cm. Démontrer qu'il n'est pas rectangle.
\end{exercice}

\corriges
\begin{corrige}[22]
Le triangle est rectangle en $A$, donc $BC$ est l'hypoténuse. Par Pythagore :
\[
BC^2=AB^2+AC^2=6^2+8^2=36+64=100.
\]
Donc $BC=10$ cm.
\end{corrige}

\begin{corrige}[23]
Le plus long côté mesure $13$ cm. On calcule :
\[
13^2=169 \quad \text{et} \quad 5^2+12^2=25+144=169.
\]
Les deux résultats sont égaux. D'après la réciproque du théorème de Pythagore, le triangle est rectangle.
\end{corrige}

\begin{corrige}[24]
Le plus long côté mesure $9$ cm. On calcule :
\[
9^2=81 \quad \text{et} \quad 4^2+7^2=16+49=65.
\]
Comme $81\neq65$, d'après la contraposée du théorème de Pythagore, le triangle n'est pas rectangle.
\end{corrige}

%====================================================
\chapitre{Théorème de Thalès}

\objectif{
\begin{itemize}
\item Reconnaître une configuration de Thalès.
\item Calculer une longueur avec Thalès.
\item Démontrer un parallélisme avec la réciproque.
\item Rédiger les rapports dans le bon ordre.
\end{itemize}}

\defs
\begin{definitionbox}
Une configuration de Thalès apparaît lorsque deux droites sécantes sont coupées par deux droites parallèles.
\end{definitionbox}

\props
\begin{proprietebox}[title=Théorème de Thalès]
Soient $A$, $B$, $M$ alignés et $A$, $C$, $N$ alignés. Si $(BC)\parallel(MN)$, alors :
\[
\frac{AM}{AB}=\frac{AN}{AC}=\frac{MN}{BC}.
\]
\end{proprietebox}

\begin{proprietebox}[title=Réciproque de Thalès]
Soient $A$, $B$, $M$ alignés dans le même ordre et $A$, $C$, $N$ alignés dans le même ordre. Si :
\[
\frac{AM}{AB}=\frac{AN}{AC},
\]
alors $(MN)\parallel(BC)$.
\end{proprietebox}

\demos
\begin{preuvebox}
La justification de Thalès repose sur l'agrandissement-réduction. Lorsque $(MN)$ est parallèle à $(BC)$, les triangles $AMN$ et $ABC$ ont la même forme : ils sont semblables. Les longueurs correspondantes sont donc proportionnelles, d'où les égalités de rapports.
\end{preuvebox}

\methodes
\begin{methodebox}
Pour utiliser Thalès :
\begin{enumerate}
\item vérifier les alignements ;
\item vérifier le parallélisme ;
\item écrire les rapports de longueurs correspondantes ;
\item remplacer par les valeurs connues ;
\item résoudre le produit en croix.
\end{enumerate}
\end{methodebox}

\exemples
\begin{exemplebox}
On sait que $A,B,M$ sont alignés, $A,C,N$ sont alignés et $(MN)\parallel(BC)$. On donne $AB=6$, $AM=3$ et $AC=10$. Alors :
\[
\frac{AM}{AB}=\frac{AN}{AC}.
\]
Donc :
\[
\frac{3}{6}=\frac{AN}{10}.
\]
Ainsi $AN=\frac{3\times10}{6}=5$.
\end{exemplebox}

\erreurs
\begin{attentionbox}
L'erreur la plus fréquente est de mélanger les rapports. Il faut toujours comparer les longueurs qui se correspondent dans les deux triangles.
\end{attentionbox}

\exos
\begin{exercice}[Calcul avec Thalès]
Dans une configuration de Thalès, on sait que $(MN)\parallel(BC)$, $AB=9$, $AM=6$ et $AC=12$. Calculer $AN$.
\end{exercice}

\exosdemo
\begin{exercice}[Réciproque]
Les points $A,B,M$ sont alignés dans le même ordre, ainsi que $A,C,N$. On sait que $AM=4$, $AB=10$, $AN=6$ et $AC=15$. Démontrer que $(MN)\parallel(BC)$.
\end{exercice}

\corriges
\begin{corrige}[25]
D'après Thalès :
\[
\frac{AM}{AB}=\frac{AN}{AC}.
\]
Donc :
\[
\frac{6}{9}=\frac{AN}{12}.
\]
Ainsi $AN=\dfrac{6\times12}{9}=8$.
\end{corrige}

\begin{corrige}[26]
On calcule :
\[
\frac{AM}{AB}=\frac{4}{10}=0,4,
\quad
\frac{AN}{AC}=\frac{6}{15}=0,4.
\]
Les rapports sont égaux. Comme les points sont alignés dans le même ordre, d'après la réciproque du théorème de Thalès, $(MN)\parallel(BC)$.
\end{corrige}

%====================================================
\chapitre{Trigonométrie dans le triangle rectangle}

\objectif{
\begin{itemize}
\item Identifier hypoténuse, côté adjacent et côté opposé.
\item Utiliser sinus, cosinus et tangente.
\item Calculer une longueur ou un angle.
\end{itemize}}

\defs
\begin{definitionbox}
Dans un triangle rectangle, relativement à un angle aigu :
\begin{itemize}
\item l'\textbf{hypoténuse} est le côté opposé à l'angle droit ;
\item le \textbf{côté opposé} est en face de l'angle étudié ;
\item le \textbf{côté adjacent} touche l'angle étudié sans être l'hypoténuse.
\end{itemize}
\end{definitionbox}

\props
\begin{proprietebox}
Dans un triangle rectangle :
\[
\cos(\theta)=\frac{\text{adjacent}}{\text{hypoténuse}},
\quad
\sin(\theta)=\frac{\text{opposé}}{\text{hypoténuse}},
\quad
\tan(\theta)=\frac{\text{opposé}}{\text{adjacent}}.
\]
\end{proprietebox}

\methodes
\begin{methodebox}
Pour choisir la bonne formule :
\begin{enumerate}
\item repérer l'angle étudié ;
\item nommer les côtés : opposé, adjacent, hypoténuse ;
\item garder seulement les deux longueurs utiles ;
\item choisir sinus, cosinus ou tangente.
\end{enumerate}
\end{methodebox}

\exemples
\begin{exemplebox}
Dans un triangle rectangle, l'hypoténuse mesure $10$ cm et l'angle $\theta$ vérifie que le côté adjacent mesure $6$ cm. Alors :
\[
\cos(\theta)=\frac{6}{10}=0,6.
\]
Donc $\theta=\arccos(0,6)\approx53,1^\circ$.
\end{exemplebox}

\erreurs
\begin{attentionbox}
Le côté opposé et le côté adjacent dépendent de l'angle choisi. Si on change d'angle, ces noms peuvent changer.
\end{attentionbox}

\exos
\begin{exercice}[Cosinus]
Dans un triangle rectangle, l'hypoténuse mesure $13$ cm et le côté adjacent à un angle $\theta$ mesure $5$ cm. Donner $\cos(\theta)$.
\end{exercice}

\exosdemo
\begin{exercice}[Choix d'une formule]
Expliquer pourquoi, si l'on connaît le côté opposé et l'hypoténuse, il faut utiliser le sinus.
\end{exercice}

\corriges
\begin{corrige}[27]
Par définition :
\[
\cos(\theta)=\frac{adjacent}{hypoténuse}=\frac{5}{13}.
\]
\end{corrige}

\begin{corrige}[28]
Dans un triangle rectangle, le sinus d'un angle aigu est défini par le rapport :
\[
\sin(\theta)=\frac{opposé}{hypoténuse}.
\]
Donc si l'on connaît le côté opposé et l'hypoténuse, le sinus est la formule adaptée.
\end{corrige}

%====================================================
\chapitre{Géométrie repérée}

\objectif{
\begin{itemize}
\item Lire et placer des coordonnées.
\item Calculer les coordonnées d'un milieu.
\item Calculer une longueur dans un repère orthonormé.
\item Utiliser les coordonnées pour démontrer.
\end{itemize}}

\defs
\begin{definitionbox}
Dans un repère du plan, un point $A$ est repéré par deux nombres appelés coordonnées :
\[
A(x_A;y_A).
\]
Le nombre $x_A$ est l'abscisse et $y_A$ est l'ordonnée.
\end{definitionbox}

\props
\begin{proprietebox}
Si $A(x_A;y_A)$ et $B(x_B;y_B)$, alors le milieu $M$ de $[AB]$ a pour coordonnées :
\[
M\left(\frac{x_A+x_B}{2};\frac{y_A+y_B}{2}\right).
\]
\end{proprietebox}

\begin{proprietebox}
Dans un repère orthonormé :
\[
AB=\sqrt{(x_B-x_A)^2+(y_B-y_A)^2}.
\]
\end{proprietebox}

\demos
\begin{preuvebox}
La formule de distance vient du théorème de Pythagore. Les points $A$ et $B$ définissent un triangle rectangle dont les côtés de l'angle droit mesurent $|x_B-x_A|$ horizontalement et $|y_B-y_A|$ verticalement. L'hypoténuse est $AB$, donc :
\[
AB^2=(x_B-x_A)^2+(y_B-y_A)^2.
\]
\end{preuvebox}

\methodes
\begin{methodebox}
Pour démontrer qu'un quadrilatère est un parallélogramme avec les coordonnées, on peut calculer les milieux de ses diagonales. Si les diagonales ont le même milieu, alors le quadrilatère est un parallélogramme.
\end{methodebox}

\exemples
\begin{exemplebox}
Soient $A(1;2)$ et $B(5;8)$. Le milieu $M$ de $[AB]$ est :
\[
M\left(\frac{1+5}{2};\frac{2+8}{2}\right)=M(3;5).
\]
\end{exemplebox}

\erreurs
\begin{attentionbox}
Pour le milieu, on moyenne les abscisses entre elles et les ordonnées entre elles. On ne mélange pas abscisse et ordonnée.
\end{attentionbox}

\exos
\begin{exercice}[Milieu]
Soient $A(2;-1)$ et $B(8;5)$. Calculer les coordonnées du milieu de $[AB]$.
\end{exercice}

\begin{exercice}[Distance]
Soient $A(1;1)$ et $B(4;5)$. Calculer $AB$.
\end{exercice}

\exosdemo
\begin{exercice}[Parallélogramme]
On donne $A(0;0)$, $B(4;1)$, $C(5;5)$ et $D(1;4)$. Démontrer que $ABCD$ est un parallélogramme.
\end{exercice}

\corriges
\begin{corrige}[29]
Le milieu $M$ a pour coordonnées :
\[
M\left(\frac{2+8}{2};\frac{-1+5}{2}\right)=M(5;2).
\]
\end{corrige}

\begin{corrige}[30]
\[
AB=\sqrt{(4-1)^2+(5-1)^2}=\sqrt{3^2+4^2}=\sqrt{25}=5.
\]
\end{corrige}

\begin{corrige}[31]
Le milieu de $[AC]$ est :
\[
M\left(\frac{0+5}{2};\frac{0+5}{2}\right)=M(2,5;2,5).
\]
Le milieu de $[BD]$ est :
\[
N\left(\frac{4+1}{2};\frac{1+4}{2}\right)=N(2,5;2,5).
\]
Les diagonales ont le même milieu. Donc $ABCD$ est un parallélogramme.
\end{corrige}

%====================================================
\chapitre{Raisonnement et démonstration en géométrie}

\objectif{
\begin{itemize}
\item Distinguer observation, conjecture et démonstration.
\item Construire une preuve claire.
\item Choisir le bon théorème.
\item Rédiger avec précision.
\end{itemize}}

\defs
\begin{definitionbox}
Une \textbf{observation} est ce que l'on voit sur une figure.

Une \textbf{conjecture} est une affirmation que l'on pense vraie après observation.

Une \textbf{démonstration} est un raisonnement qui prouve qu'une affirmation est vraie à partir de données et de propriétés connues.
\end{definitionbox}

\props
\begin{proprietebox}
Une démonstration doit contenir :
\begin{enumerate}
\item les données utilisées ;
\item la propriété ou le théorème appliqué ;
\item la conclusion.
\end{enumerate}
\end{proprietebox}

\methodes
\begin{methodebox}
Pour choisir le bon théorème :
\begin{itemize}
\item longueur dans triangle rectangle : Pythagore ;
\item prouver triangle rectangle : réciproque de Pythagore ;
\item longueurs proportionnelles avec parallèles : Thalès ;
\item prouver parallèles avec rapports : réciproque de Thalès ;
\item égalités de distances : médiatrice, cercle ou symétrie ;
\item égalités d'angles : triangles particuliers, angles opposés, angles alternes-internes.
\end{itemize}
\end{methodebox}

\exemples
\begin{exemplebox}
Démontrer que $ABC$ est rectangle sachant que $AB=6$, $AC=8$ et $BC=10$.

Le plus long côté est $BC$. On calcule :
\[
BC^2=10^2=100,
\quad AB^2+AC^2=6^2+8^2=36+64=100.
\]
Donc $BC^2=AB^2+AC^2$. D'après la réciproque du théorème de Pythagore, le triangle $ABC$ est rectangle en $A$.
\end{exemplebox}

\erreurs
\begin{attentionbox}
Écrire seulement « d'après Pythagore » ne suffit pas. Il faut préciser dans quel triangle, pourquoi on peut appliquer le théorème, puis écrire l'égalité correcte.
\end{attentionbox}

\exos
\begin{exercice}[Rédaction]
Réécrire proprement la démonstration suivante : « C'est rectangle parce que $3^2+4^2=5^2$. »
\end{exercice}

\exosdemo
\begin{exercice}[Choix du théorème]
Pour chaque objectif, indiquer un théorème possible :
\begin{enumerate}
\item prouver qu'un triangle est rectangle ;
\item calculer une longueur dans une configuration avec deux parallèles ;
\item prouver qu'un point appartient à une médiatrice.
\end{enumerate}
\end{exercice}

\corriges
\begin{corrige}[32]
Dans le triangle considéré, le plus long côté mesure $5$. On calcule :
\[
5^2=25 \quad \text{et} \quad 3^2+4^2=9+16=25.
\]
Donc $5^2=3^2+4^2$. D'après la réciproque du théorème de Pythagore, le triangle est rectangle.
\end{corrige}

\begin{corrige}[33]
\begin{enumerate}
\item Pour prouver qu'un triangle est rectangle : réciproque du théorème de Pythagore.
\item Pour calculer une longueur avec deux parallèles : théorème de Thalès.
\item Pour prouver qu'un point appartient à une médiatrice : propriété caractéristique de la médiatrice, en montrant qu'il est équidistant des deux extrémités du segment.
\end{enumerate}
\end{corrige}

%====================================================
\clearpage
\section{Exercices bilan par grands thèmes}

\begin{exercice}[Bilan angles et parallèles]
Deux droites parallèles sont coupées par une sécante. Un angle correspondant mesure $128^\circ$.
\begin{enumerate}
\item Donner la mesure de l'angle correspondant associé.
\item Donner la mesure de l'angle supplémentaire.
\item Rédiger une justification complète.
\end{enumerate}
\end{exercice}

\begin{exercice}[Bilan triangles et quadrilatères]
$ABCD$ est un quadrilatère dont les diagonales se coupent en $O$. On sait que $OA=OC$, $OB=OD$ et $AB=BC$.
\begin{enumerate}
\item Démontrer que $ABCD$ est un parallélogramme.
\item Démontrer que $ABCD$ est un losange.
\end{enumerate}
\end{exercice}

\begin{exercice}[Bilan Pythagore-Thalès]
Dans un triangle $ABC$, les points $M$ et $N$ appartiennent respectivement à $[AB]$ et $[AC]$. On sait que $(MN)\parallel(BC)$, $AM=4$, $AB=10$, $AC=15$.
\begin{enumerate}
\item Calculer $AN$.
\item Si $MN=6$, calculer $BC$.
\end{enumerate}
\end{exercice}

\begin{exercice}[Bilan repère]
Dans un repère orthonormé, on donne $A(0;0)$, $B(6;0)$, $C(6;8)$.
\begin{enumerate}
\item Calculer $AB$, $BC$ et $AC$.
\item Démontrer que le triangle $ABC$ est rectangle.
\item Calculer son aire.
\end{enumerate}
\end{exercice}

%====================================================
\clearpage
\section{Grand devoir bilan final}

\subsection*{Exercice 1 -- Angles et parallélisme}
Deux droites $(d)$ et $(d')$ sont coupées par une sécante $(\Delta)$. Deux angles alternes-internes mesurent $63^\circ$.
\begin{enumerate}
\item Que peut-on dire des droites $(d)$ et $(d')$ ? Justifier.
\item Un angle adjacent à l'un de ces angles est noté $x$. Calculer $x$.
\end{enumerate}

\subsection*{Exercice 2 -- Triangle et cercle}
Soit $ABC$ un triangle isocèle en $A$ tel que $\widehat{B}=48^\circ$.
\begin{enumerate}
\item Calculer $\widehat{C}$ et $\widehat{A}$.
\item On suppose que $O$ est équidistant de $A$, $B$ et $C$. Que représente $O$ pour le triangle $ABC$ ?
\end{enumerate}

\subsection*{Exercice 3 -- Pythagore}
Un triangle $RST$ vérifie $RS=9$ cm, $RT=12$ cm et $ST=15$ cm.
\begin{enumerate}
\item Montrer que le triangle est rectangle.
\item Calculer son aire.
\end{enumerate}

\subsection*{Exercice 4 -- Thalès}
Dans une configuration de Thalès, $A,B,M$ sont alignés, $A,C,N$ sont alignés et $(MN)\parallel(BC)$. On donne $AM=5$, $AB=8$, $AC=12$ et $BC=10$.
\begin{enumerate}
\item Calculer $AN$.
\item Calculer $MN$.
\end{enumerate}

\subsection*{Exercice 5 -- Géométrie repérée}
On donne $A(1;2)$, $B(7;2)$, $C(7;10)$ et $D(1;10)$.
\begin{enumerate}
\item Calculer les longueurs $AB$ et $BC$.
\item Démontrer que $ABCD$ est un rectangle.
\item Calculer son aire.
\end{enumerate}

%====================================================
\clearpage
\section{Correction détaillée du grand devoir bilan}

\subsection*{Correction de l'exercice 1}
Deux angles alternes-internes sont égaux. Or si deux angles alternes-internes formés par deux droites et une sécante sont égaux, alors les deux droites sont parallèles. Donc $(d)\parallel(d')$.

L'angle adjacent à un angle de $63^\circ$ forme avec lui un angle plat. Donc :
\[
x=180^\circ-63^\circ=117^\circ.
\]

\subsection*{Correction de l'exercice 2}
Le triangle $ABC$ est isocèle en $A$, donc les angles à la base sont égaux :
\[
\widehat{B}=\widehat{C}=48^\circ.
\]
La somme des angles d'un triangle vaut $180^\circ$, donc :
\[
\widehat{A}=180^\circ-48^\circ-48^\circ=84^\circ.
\]
Si $O$ est équidistant de $A$, $B$ et $C$, alors $O$ est le centre du cercle passant par les trois sommets du triangle. C'est donc le centre du cercle circonscrit au triangle $ABC$.

\subsection*{Correction de l'exercice 3}
Le plus long côté est $ST=15$ cm. On calcule :
\[
ST^2=15^2=225,
\quad RS^2+RT^2=9^2+12^2=81+144=225.
\]
Donc $ST^2=RS^2+RT^2$. D'après la réciproque du théorème de Pythagore, le triangle $RST$ est rectangle en $R$.

Son aire vaut :
\[
\mathcal A=\frac{RS\times RT}{2}=\frac{9\times12}{2}=54\text{ cm}^2.
\]

\subsection*{Correction de l'exercice 4}
D'après le théorème de Thalès :
\[
\frac{AM}{AB}=\frac{AN}{AC}=\frac{MN}{BC}.
\]
Donc :
\[
\frac{5}{8}=\frac{AN}{12}.
\]
Ainsi :
\[
AN=\frac{5\times12}{8}=7,5.
\]
Puis :
\[
\frac{5}{8}=\frac{MN}{10},
\]
donc :
\[
MN=\frac{5\times10}{8}=6,25.
\]

\subsection*{Correction de l'exercice 5}
Les points $A$ et $B$ ont la même ordonnée, donc $AB=7-1=6$. Les points $B$ et $C$ ont la même abscisse, donc $BC=10-2=8$.

On observe que $AB$ est horizontal et $BC$ est vertical ; ainsi $(AB)\perp(BC)$. De plus, $AD$ est vertical et $BC$ est vertical, donc $(AD)\parallel(BC)$. De même, $AB$ et $DC$ sont horizontales, donc $(AB)\parallel(DC)$. Le quadrilatère possède deux paires de côtés opposés parallèles et un angle droit : c'est un rectangle.

Son aire vaut :
\[
\mathcal A=AB\times BC=6\times8=48.
\]

%====================================================
\clearpage
\section{Fiche de synthèse finale}

\begin{retenirbox}
\textbf{Milieu :} $M$ est le milieu de $[AB]$ si $M\in[AB]$ et $AM=MB$.
\end{retenirbox}

\begin{retenirbox}
\textbf{Angles :} deux angles opposés par le sommet sont égaux. Deux angles supplémentaires ont une somme de $180^\circ$.
\end{retenirbox}

\begin{retenirbox}
\textbf{Parallèles :} si deux droites sont perpendiculaires à une même troisième droite, alors elles sont parallèles. Si deux droites parallèles sont coupées par une sécante, alors les angles alternes-internes et correspondants sont égaux.
\end{retenirbox}

\begin{retenirbox}
\textbf{Triangles :} la somme des angles d'un triangle vaut $180^\circ$. Dans un triangle isocèle, les angles à la base sont égaux. Dans un triangle équilatéral, les trois angles valent $60^\circ$.
\end{retenirbox}

\begin{retenirbox}
\textbf{Quadrilatères :} un parallélogramme a ses côtés opposés parallèles et ses diagonales qui se coupent en leur milieu. Un rectangle a quatre angles droits. Un losange a quatre côtés égaux. Un carré est à la fois rectangle et losange.
\end{retenirbox}

\begin{retenirbox}
\textbf{Médiatrice :} un point appartient à la médiatrice de $[AB]$ si et seulement s'il est équidistant de $A$ et de $B$.
\end{retenirbox}

\begin{retenirbox}
\textbf{Pythagore :} si $ABC$ est rectangle en $A$, alors $BC^2=AB^2+AC^2$. La réciproque permet de prouver qu'un triangle est rectangle.
\end{retenirbox}

\begin{retenirbox}
\textbf{Thalès :} dans une configuration avec deux droites sécantes coupées par deux parallèles, les longueurs correspondantes sont proportionnelles.
\end{retenirbox}

\begin{retenirbox}
\textbf{Trigonométrie :}
\[
\cos=\frac{adjacent}{hypoténuse},
\quad
\sin=\frac{opposé}{hypoténuse},
\quad
\tan=\frac{opposé}{adjacent}.
\]
\end{retenirbox}

\begin{retenirbox}
\textbf{Coordonnées :}
\[
M\left(\frac{x_A+x_B}{2};\frac{y_A+y_B}{2}\right),
\quad
AB=\sqrt{(x_B-x_A)^2+(y_B-y_A)^2}.
\]
\end{retenirbox}

\begin{center}
\Large\bfseries Fin du cours
\end{center}


%====================================================
\clearpage
\section{Banque supplémentaire d'exercices progressifs}

Cette banque d'exercices peut être utilisée en devoir maison, en entraînement différencié ou en préparation au brevet. Les exercices sont classés par grands thèmes.

%====================================================
\subsection{Thème 1 -- Bases, alignement, milieux et longueurs}

\begin{exercice}[Calculs de longueurs]
Les points $A$, $B$ et $C$ sont alignés dans cet ordre. On donne $AB=4,8$ cm et $BC=3,7$ cm.
\begin{enumerate}
\item Faire une figure à main levée codée.
\item Calculer $AC$.
\item Rédiger une phrase de justification.
\end{enumerate}
\end{exercice}

\begin{exercice}[Milieu]
Le point $M$ est le milieu du segment $[AB]$. On sait que $AM=7,5$ cm.
\begin{enumerate}
\item Calculer $MB$.
\item Calculer $AB$.
\item Expliquer pourquoi un point $P$ tel que $PA=PB$ n'est pas forcément le milieu de $[AB]$.
\end{enumerate}
\end{exercice}

\begin{exercice}[Reconnaître un milieu]
Les points $A$, $M$ et $B$ sont alignés dans cet ordre. On donne $AB=18$ cm et $AM=9$ cm.
\begin{enumerate}
\item Calculer $MB$.
\item Démontrer que $M$ est le milieu de $[AB]$.
\end{enumerate}
\end{exercice}

%====================================================
\subsection{Thème 2 -- Angles et droites parallèles}

\begin{exercice}[Angles supplémentaires]
Deux angles adjacents forment un angle plat. L'un mesure $37^\circ$.
\begin{enumerate}
\item Calculer la mesure de l'autre angle.
\item Dire si les deux angles sont supplémentaires.
\end{enumerate}
\end{exercice}

\begin{exercice}[Angles opposés par le sommet]
Deux droites $(AB)$ et $(CD)$ se coupent en $O$. On sait que $\widehat{AOC}=112^\circ$.
\begin{enumerate}
\item Donner la mesure de l'angle opposé par le sommet à $\widehat{AOC}$.
\item Calculer la mesure de l'angle adjacent à $\widehat{AOC}$.
\end{enumerate}
\end{exercice}

\begin{exercice}[Parallélisme par les angles]
Deux droites $(d_1)$ et $(d_2)$ sont coupées par une sécante. Deux angles alternes-internes mesurent chacun $58^\circ$.
\begin{enumerate}
\item Que peut-on conclure sur les droites $(d_1)$ et $(d_2)$ ?
\item Rédiger une démonstration complète.
\end{enumerate}
\end{exercice}

%====================================================
\subsection{Thème 3 -- Triangles particuliers}

\begin{exercice}[Triangle isocèle]
Le triangle $ABC$ est isocèle en $A$. On sait que $\widehat{A}=40^\circ$.
\begin{enumerate}
\item Quels sont les deux angles égaux ?
\item Calculer $\widehat{B}$ et $\widehat{C}$.
\end{enumerate}
\end{exercice}

\begin{exercice}[Triangle équilatéral]
Le triangle $DEF$ est équilatéral de côté $6$ cm.
\begin{enumerate}
\item Donner la mesure de chacun de ses angles.
\item Calculer son périmètre.
\end{enumerate}
\end{exercice}

\begin{exercice}[Inégalité triangulaire]
Peut-on construire un triangle dont les côtés mesurent :
\begin{enumerate}
\item $4$ cm, $5$ cm et $8$ cm ?
\item $3$ cm, $6$ cm et $10$ cm ?
\item $7$ cm, $7$ cm et $14$ cm ?
\end{enumerate}
Justifier chaque réponse.
\end{exercice}

%====================================================
\subsection{Thème 4 -- Quadrilatères}

\begin{exercice}[Reconnaissance d'un parallélogramme]
Dans un quadrilatère $ABCD$, les diagonales $[AC]$ et $[BD]$ se coupent en $O$. On sait que $OA=OC$ et $OB=OD$.
\begin{enumerate}
\item Que représente $O$ pour $[AC]$ ?
\item Que représente $O$ pour $[BD]$ ?
\item Démontrer que $ABCD$ est un parallélogramme.
\end{enumerate}
\end{exercice}

\begin{exercice}[Rectangle, losange ou carré]
$ABCD$ est un parallélogramme.
\begin{enumerate}
\item On sait que $AB=BC$. Que peut-on conclure ?
\item On sait que $\widehat{ABC}=90^\circ$. Que peut-on conclure ?
\item Si les deux informations sont vraies simultanément, que peut-on conclure ?
\end{enumerate}
\end{exercice}

%====================================================
\subsection{Thème 5 -- Cercles et médiatrices}

\begin{exercice}[Vocabulaire du cercle]
Un cercle a pour centre $O$ et pour rayon $5$ cm. Le point $A$ appartient au cercle.
\begin{enumerate}
\item Quelle est la longueur $OA$ ?
\item Quelle est la longueur d'un diamètre du cercle ?
\item Si $B$ appartient aussi au cercle, peut-on dire que $OA=OB$ ? Justifier.
\end{enumerate}
\end{exercice}

\begin{exercice}[Médiatrice]
On donne deux points $A$ et $B$. Le point $P$ vérifie $PA=PB$.
\begin{enumerate}
\item Quelle propriété peut-on utiliser ?
\item Que peut-on conclure sur le point $P$ ?
\end{enumerate}
\end{exercice}

%====================================================
\subsection{Thème 6 -- Transformations}

\begin{exercice}[Symétrie axiale]
Le point $A'$ est l'image de $A$ par la symétrie d'axe $(d)$.
\begin{enumerate}
\item Que représente la droite $(d)$ pour le segment $[AA']$ ?
\item Quelle égalité de distances peut-on écrire si $M$ est un point de l'axe $(d)$ ?
\end{enumerate}
\end{exercice}

\begin{exercice}[Symétrie centrale]
Le point $B'$ est l'image de $B$ par la symétrie centrale de centre $O$.
\begin{enumerate}
\item Que représente $O$ pour $[BB']$ ?
\item Si $OB=4,2$ cm, calculer $BB'$.
\end{enumerate}
\end{exercice}

\begin{exercice}[Homothétie]
Une homothétie de rapport $3$ transforme un segment $[AB]$ de longueur $4$ cm en un segment $[A'B']$.
\begin{enumerate}
\item Calculer $A'B'$.
\item Par combien l'aire d'une figure est-elle multipliée ?
\end{enumerate}
\end{exercice}

%====================================================
\subsection{Thème 7 -- Aires, périmètres et volumes}

\begin{exercice}[Aire composée]
Une figure est composée d'un rectangle de longueur $12$ cm et de largeur $5$ cm, auquel on ajoute un triangle rectangle de base $5$ cm et de hauteur $4$ cm.
\begin{enumerate}
\item Calculer l'aire du rectangle.
\item Calculer l'aire du triangle.
\item Calculer l'aire totale de la figure.
\end{enumerate}
\end{exercice}

\begin{exercice}[Cercle et disque]
Un disque a pour rayon $6$ cm.
\begin{enumerate}
\item Calculer son aire exacte.
\item Calculer son périmètre exact.
\item Donner une valeur approchée de l'aire au dixième près.
\end{enumerate}
\end{exercice}

\begin{exercice}[Volume d'un cylindre]
Un cylindre a pour rayon $4$ cm et pour hauteur $9$ cm.
\begin{enumerate}
\item Calculer l'aire de sa base.
\item Calculer son volume exact.
\end{enumerate}
\end{exercice}

\begin{exercice}[Volume d'une pyramide]
Une pyramide a pour base un carré de côté $6$ cm et pour hauteur $10$ cm.
\begin{enumerate}
\item Calculer l'aire de la base.
\item Calculer le volume de la pyramide.
\end{enumerate}
\end{exercice}

%====================================================
\subsection{Thème 8 -- Pythagore, Thalès et trigonométrie}

\begin{exercice}[Pythagore direct]
Un triangle $ABC$ est rectangle en $A$. On donne $AB=7$ cm et $AC=24$ cm.
\begin{enumerate}
\item Écrire l'égalité de Pythagore.
\item Calculer $BC$.
\end{enumerate}
\end{exercice}

\begin{exercice}[Pythagore réciproque]
Un triangle $DEF$ a pour côtés $DE=8$ cm, $DF=15$ cm et $EF=17$ cm.
\begin{enumerate}
\item Identifier le plus long côté.
\item Vérifier si l'égalité de Pythagore est vraie.
\item Conclure.
\end{enumerate}
\end{exercice}

\begin{exercice}[Thalès direct]
Dans un triangle $ABC$, les points $M$ et $N$ appartiennent respectivement aux côtés $[AB]$ et $[AC]$. On sait que $(MN)\parallel(BC)$, $AB=12$, $AM=8$ et $AC=15$.
\begin{enumerate}
\item Écrire les rapports de Thalès.
\item Calculer $AN$.
\end{enumerate}
\end{exercice}

\begin{exercice}[Réciproque de Thalès]
Les points $A,M,B$ sont alignés dans cet ordre et les points $A,N,C$ sont alignés dans cet ordre. On donne $AM=6$, $AB=9$, $AN=10$ et $AC=15$.
\begin{enumerate}
\item Calculer $\dfrac{AM}{AB}$ et $\dfrac{AN}{AC}$.
\item Que peut-on conclure sur les droites $(MN)$ et $(BC)$ ?
\end{enumerate}
\end{exercice}

\begin{exercice}[Trigonométrie]
Dans un triangle rectangle, l'hypoténuse mesure $20$ cm et un angle aigu mesure $35^\circ$.
\begin{enumerate}
\item Calculer la longueur du côté adjacent à cet angle.
\item Calculer la longueur du côté opposé à cet angle.
\end{enumerate}
\end{exercice}

%====================================================
\subsection{Thème 9 -- Géométrie repérée}

\begin{exercice}[Coordonnées du milieu]
On donne $A(-2;4)$ et $B(6;10)$.
\begin{enumerate}
\item Calculer les coordonnées du milieu $M$ de $[AB]$.
\item Vérifier que les abscisses et les ordonnées ont bien été traitées séparément.
\end{enumerate}
\end{exercice}

\begin{exercice}[Distance dans un repère]
Dans un repère orthonormé, on donne $C(-1;2)$ et $D(5;10)$.
\begin{enumerate}
\item Calculer $CD$.
\item Interpréter ce calcul comme une application du théorème de Pythagore.
\end{enumerate}
\end{exercice}

\begin{exercice}[Nature d'un triangle]
On donne $A(0;0)$, $B(5;0)$ et $C(0;12)$.
\begin{enumerate}
\item Calculer $AB$, $AC$ et $BC$.
\item Démontrer que le triangle $ABC$ est rectangle.
\item Calculer son aire.
\end{enumerate}
\end{exercice}

%====================================================
\clearpage
\section{Corrigés de la banque supplémentaire}

\subsection*{Corrigés du thème 1}

\begin{corrige}[34]
Les points sont alignés dans l'ordre $A$, $B$, $C$, donc :
\[
AC=AB+BC.
\]
Ainsi :
\[
AC=4,8+3,7=8,5\text{ cm}.
\]
\end{corrige}

\begin{corrige}[35]
Puisque $M$ est le milieu de $[AB]$, on a :
\[
AM=MB.
\]
Donc :
\[
MB=7,5\text{ cm}.
\]
De plus :
\[
AB=AM+MB=7,5+7,5=15\text{ cm}.
\]
Un point $P$ tel que $PA=PB$ est équidistant de $A$ et de $B$. Il appartient donc à la médiatrice de $[AB]$, mais il n'est pas forcément situé sur le segment $[AB]$. Il n'est donc pas forcément le milieu de $[AB]$.
\end{corrige}

\begin{corrige}[36]
Comme les points $A$, $M$ et $B$ sont alignés dans cet ordre :
\[
MB=AB-AM=18-9=9\text{ cm}.
\]
On a donc $AM=MB$ et $M\in[AB]$. Par définition, $M$ est le milieu de $[AB]$.
\end{corrige}

\subsection*{Corrigés du thème 2}

\begin{corrige}[37]
Deux angles formant un angle plat ont une somme égale à $180^\circ$. L'autre angle mesure donc :
\[
180^\circ-37^\circ=143^\circ.
\]
Les deux angles sont supplémentaires car leur somme vaut $180^\circ$.
\end{corrige}

\begin{corrige}[38]
Deux angles opposés par le sommet ont la même mesure. L'angle opposé à $\widehat{AOC}$ mesure donc $112^\circ$.

L'angle adjacent à $\widehat{AOC}$ est supplémentaire avec lui, donc il mesure :
\[
180^\circ-112^\circ=68^\circ.
\]
\end{corrige}

\begin{corrige}[39]
Les deux angles alternes-internes sont égaux. Or, si deux angles alternes-internes formés par deux droites et une sécante sont égaux, alors ces deux droites sont parallèles. Donc :
\[
(d_1)\parallel(d_2).
\]
\end{corrige}

\subsection*{Corrigés du thème 3}

\begin{corrige}[40]
Le triangle $ABC$ est isocèle en $A$, donc les angles à la base sont égaux :
\[
\widehat{B}=\widehat{C}.
\]
Comme la somme des angles d'un triangle vaut $180^\circ$ :
\[
\widehat{B}+\widehat{C}=180^\circ-40^\circ=140^\circ.
\]
Donc :
\[
\widehat{B}=\widehat{C}=70^\circ.
\]
\end{corrige}

\begin{corrige}[41]
Dans un triangle équilatéral, les trois angles mesurent $60^\circ$.

Le périmètre du triangle vaut :
\[
P=6+6+6=18\text{ cm}.
\]
\end{corrige}

\begin{corrige}[42]
On vérifie l'inégalité triangulaire.

\begin{enumerate}
\item Pour $4$, $5$ et $8$ : le plus grand côté est $8$ et $8<4+5=9$. Le triangle est constructible.
\item Pour $3$, $6$ et $10$ : le plus grand côté est $10$ et $10>3+6=9$. Le triangle n'est pas constructible.
\item Pour $7$, $7$ et $14$ : le plus grand côté est $14$ et $14=7+7$. On obtient un triangle aplati, donc ce n'est pas un vrai triangle.
\end{enumerate}
\end{corrige}

\subsection*{Corrigés du thème 4}

\begin{corrige}[43]
Puisque $OA=OC$, le point $O$ est le milieu de $[AC]$.

Puisque $OB=OD$, le point $O$ est le milieu de $[BD]$.

Les diagonales du quadrilatère $ABCD$ se coupent donc en leur milieu. Or, si les diagonales d'un quadrilatère se coupent en leur milieu, alors ce quadrilatère est un parallélogramme. Donc $ABCD$ est un parallélogramme.
\end{corrige}

\begin{corrige}[44]
\begin{enumerate}
\item Si $ABCD$ est un parallélogramme et si deux côtés consécutifs sont égaux, alors $ABCD$ est un losange.
\item Si $ABCD$ est un parallélogramme et s'il possède un angle droit, alors $ABCD$ est un rectangle.
\item Si les deux informations sont vraies, alors $ABCD$ est à la fois un rectangle et un losange. C'est donc un carré.
\end{enumerate}
\end{corrige}

\subsection*{Corrigés du thème 5}

\begin{corrige}[45]
Si $A$ appartient au cercle de centre $O$ et de rayon $5$ cm, alors :
\[
OA=5\text{ cm}.
\]
Le diamètre vaut deux fois le rayon :
\[
d=2\times5=10\text{ cm}.
\]
Si $B$ appartient aussi au cercle, alors $OB=5$ cm. Donc :
\[
OA=OB.
\]
\end{corrige}

\begin{corrige}[46]
On sait que :
\[
PA=PB.
\]
Donc $P$ est équidistant de $A$ et de $B$. Or l'ensemble des points équidistants de $A$ et de $B$ est la médiatrice de $[AB]$. Donc $P$ appartient à la médiatrice de $[AB]$.
\end{corrige}

\subsection*{Corrigés du thème 6}

\begin{corrige}[47]
Dans une symétrie axiale, l'axe est la médiatrice du segment reliant un point à son image. Donc la droite $(d)$ est la médiatrice de $[AA']$.

Si $M$ est un point de l'axe $(d)$, alors $M$ appartient à la médiatrice de $[AA']$, donc :
\[
MA=MA'.
\]
\end{corrige}

\begin{corrige}[48]
Dans une symétrie centrale de centre $O$, le centre est le milieu du segment reliant un point et son image. Donc $O$ est le milieu de $[BB']$.

Ainsi :
\[
BB'=2\times OB=2\times4,2=8,4\text{ cm}.
\]
\end{corrige}

\begin{corrige}[49]
Une homothétie de rapport $3$ multiplie les longueurs par $3$. Donc :
\[
A'B'=3\times4=12\text{ cm}.
\]
Les aires sont multipliées par le carré du rapport :
\[
3^2=9.
\]
L'aire est donc multipliée par $9$.
\end{corrige}

\subsection*{Corrigés du thème 7}

\begin{corrige}[50]
L'aire du rectangle vaut :
\[
12\times5=60\text{ cm}^2.
\]
L'aire du triangle vaut :
\[
\frac{5\times4}{2}=10\text{ cm}^2.
\]
L'aire totale vaut :
\[
60+10=70\text{ cm}^2.
\]
\end{corrige}

\begin{corrige}[51]
L'aire exacte du disque vaut :
\[
\pi r^2=\pi\times6^2=36\pi\text{ cm}^2.
\]
Le périmètre exact du cercle vaut :
\[
2\pi r=2\pi\times6=12\pi\text{ cm}.
\]
Une valeur approchée de l'aire est :
\[
36\pi\approx113,1\text{ cm}^2.
\]
\end{corrige}

\begin{corrige}[52]
L'aire de la base vaut :
\[
\pi r^2=\pi\times4^2=16\pi\text{ cm}^2.
\]
Le volume du cylindre vaut :
\[
V=16\pi\times9=144\pi\text{ cm}^3.
\]
\end{corrige}

\begin{corrige}[53]
La base est un carré de côté $6$ cm, donc son aire vaut :
\[
6^2=36\text{ cm}^2.
\]
Le volume de la pyramide vaut :
\[
V=\frac{36\times10}{3}=120\text{ cm}^3.
\]
\end{corrige}

\subsection*{Corrigés du thème 8}

\begin{corrige}[54]
Le triangle est rectangle en $A$, donc $BC$ est l'hypoténuse. D'après le théorème de Pythagore :
\[
BC^2=AB^2+AC^2.
\]
Donc :
\[
BC^2=7^2+24^2=49+576=625.
\]
Ainsi :
\[
BC=\sqrt{625}=25\text{ cm}.
\]
\end{corrige}

\begin{corrige}[55]
Le plus long côté est $EF=17$ cm. On calcule :
\[
EF^2=17^2=289,
\]
et :
\[
DE^2+DF^2=8^2+15^2=64+225=289.
\]
Donc :
\[
EF^2=DE^2+DF^2.
\]
D'après la réciproque du théorème de Pythagore, le triangle $DEF$ est rectangle en $D$.
\end{corrige}

\begin{corrige}[56]
D'après le théorème de Thalès :
\[
\frac{AM}{AB}=\frac{AN}{AC}=\frac{MN}{BC}.
\]
Donc :
\[
\frac{8}{12}=\frac{AN}{15}.
\]
Par produit en croix :
\[
AN=\frac{8\times15}{12}=10.
\]
\end{corrige}

\begin{corrige}[57]
On calcule :
\[
\frac{AM}{AB}=\frac{6}{9}=\frac{2}{3},
\]
et :
\[
\frac{AN}{AC}=\frac{10}{15}=\frac{2}{3}.
\]
Les rapports sont égaux et les points sont alignés dans le même ordre. D'après la réciproque du théorème de Thalès :
\[
(MN)\parallel(BC).
\]
\end{corrige}

\begin{corrige}[58]
Le côté adjacent à l'angle de $35^\circ$ vérifie :
\[
\cos(35^\circ)=\frac{\text{adjacent}}{20}.
\]
Donc :
\[
\text{adjacent}=20\cos(35^\circ)\approx16,4\text{ cm}.
\]

Le côté opposé vérifie :
\[
\sin(35^\circ)=\frac{\text{opposé}}{20}.
\]
Donc :
\[
\text{opposé}=20\sin(35^\circ)\approx11,5\text{ cm}.
\]
\end{corrige}

\subsection*{Corrigés du thème 9}

\begin{corrige}[59]
Le milieu $M$ de $[AB]$ a pour coordonnées :
\[
M\left(\frac{x_A+x_B}{2};\frac{y_A+y_B}{2}\right).
\]
Donc :
\[
M\left(\frac{-2+6}{2};\frac{4+10}{2}\right)
=
M(2;7).
\]
On a bien moyenné les abscisses ensemble et les ordonnées ensemble.
\end{corrige}

\begin{corrige}[60]
Dans un repère orthonormé :
\[
CD=\sqrt{(x_D-x_C)^2+(y_D-y_C)^2}.
\]
Donc :
\[
CD=\sqrt{(5-(-1))^2+(10-2)^2}
=\sqrt{6^2+8^2}
=\sqrt{36+64}
=\sqrt{100}
=10.
\]
Ce calcul vient du théorème de Pythagore appliqué au triangle rectangle formé par le déplacement horizontal et le déplacement vertical entre $C$ et $D$.
\end{corrige}

\begin{corrige}[61]
On calcule :
\[
AB=5,
\qquad
AC=12.
\]
Puis :
\[
BC=\sqrt{(5-0)^2+(0-12)^2}
=\sqrt{25+144}
=\sqrt{169}
=13.
\]
Le plus long côté est $BC=13$. On vérifie :
\[
BC^2=13^2=169,
\]
et :
\[
AB^2+AC^2=5^2+12^2=25+144=169.
\]
Donc :
\[
BC^2=AB^2+AC^2.
\]
D'après la réciproque du théorème de Pythagore, le triangle $ABC$ est rectangle en $A$.

Son aire vaut :
\[
\mathcal A=\frac{AB\times AC}{2}
=\frac{5\times12}{2}
=30.
\]
\end{corrige}

%====================================================
\clearpage
\section{Conseils de rédaction pour les élèves}

\begin{retenirbox}
Une bonne démonstration géométrique doit être lisible même sans regarder la figure. La figure aide à chercher, mais la preuve doit reposer sur les données et les propriétés.
\end{retenirbox}

\begin{methodebox}
Pour rédiger une preuve :
\begin{enumerate}
\item Je cite les données utiles.
\item Je nomme précisément la propriété ou le théorème utilisé.
\item J'applique cette propriété à la figure.
\item Je conclus clairement.
\end{enumerate}
\end{methodebox}

\begin{exemplebox}
Mauvaise rédaction :

\medskip

\og Le triangle est rectangle car ça se voit et $3^2+4^2=5^2$. \fg

\medskip

Bonne rédaction :

\medskip

Dans le triangle considéré, le plus long côté mesure $5$ cm. On calcule :
\[
5^2=25
\]
et :
\[
3^2+4^2=9+16=25.
\]
Donc :
\[
5^2=3^2+4^2.
\]
D'après la réciproque du théorème de Pythagore, le triangle est rectangle.
\end{exemplebox}

%====================================================
\clearpage
\section{Conclusion générale}

La géométrie du collège repose sur quelques idées fondamentales :

\begin{itemize}
\item savoir nommer précisément les objets ;
\item savoir lire une figure sans confondre observation et preuve ;
\item savoir utiliser les propriétés des figures usuelles ;
\item savoir choisir le bon théorème ;
\item savoir rédiger une démonstration complète.
\end{itemize}

\begin{retenirbox}
La phrase essentielle à retenir est :

\[
\text{Une figure suggère, mais seule une démonstration prouve.}
\]
\end{retenirbox}

\begin{center}
\Large\bfseries Fin du cours
\end{center}

\end{document}